% Generated draft: requires mathematical and source review.
\documentclass[11pt]{article}
\usepackage[margin=25mm]{geometry}
\usepackage{fontspec}
\setmainfont{Latin Modern Roman}
\usepackage{amsmath,amssymb,mathtools,unicode-math}
\setmathfont{Latin Modern Math}
\usepackage{enumitem,needspace,xurl,hyperref}
\setlength{\parindent}{0pt}
\setlength{\parskip}{5pt}
\setlength{\emergencystretch}{3em}
\newcommand{\N}{\mathbb{N}}
\newcommand{\Z}{\mathbb{Z}}
\newcommand{\Q}{\mathbb{Q}}
\newcommand{\R}{\mathbb{R}}
\newcommand{\C}{\mathbb{C}}
\newcommand{\F}{\mathbb{F}}
\newcommand{\A}{\mathbb{A}}
\newcommand{\PP}{\mathbb{P}}
\newcommand{\E}{\mathbb{E}}
\newcommand{\eps}{\varepsilon}
\newcommand{\dd}{\,\mathrm{d}}
\begin{document}
\section*{UnsolvedMath ID: 30004257}
\textbf{Source problem code:} catalogue entry
\subsection{Definitions and mathematical statement}
Let \(G\) be a reductive \(p\)-adic group in the scope intended by the source. Write \(\mathcal H_I\) for its Iwahori--Hecke algebra and \(c_{\mathcal H_I}\) for the operation called the \(c\)-operation in the supplied record. These are placeholders for source-specific constructions: the record does not define this operation or specify the coefficient ring, algebra conventions, or hypotheses on \(G\).

\textbf{Requested construction.} Construct a realization of the \(c\)-operation for \(G\) itself that recovers \(c_{\mathcal H_I}\) on its Iwahori--Hecke algebra.

The supplied information does not specify what mathematical object constitutes a realization for \(G\), or the comparison procedure by which it must recover the algebra-level operation. Consequently, this is a statement schema requiring clarification from the cited report, rather than a fully specified construction problem. In particular, no particular involution, representation-theoretic functor, or compatibility identity is asserted here. The supplied assessment mentions Bernstein projectives as a proposed route; it does not supply a construction or a theorem establishing the requested realization.
\subsection{Short English statement}
The problem asks for a realization of the \(c\)-operation for a reductive \(p\)-adic group that recovers the known Iwahori--Hecke algebra operation. The operation and the required comparison are undefined in the supplied record, so the target remains incomplete. The supplied open status requires independent verification.
\subsection{Sources}
\begin{itemize}
\item Supplied local UnsolvedMath registry record, permanent ID \(30004257\), title: \emph{Group-Level Realization of the c-Operation}; source problem code: \texttt{catalogue entry}. The record belongs to the Oberwolfach Reports Open Problems collection and lists proposed year \(2019\); no individual proposer is identified.
\item Supplied report citation: \emph{New Developments in Representation Theory of p-adic Groups} (2020), Oberwolfach Reports, pp.\ 2739--2819, DOI \texttt{10.4171/owr/2019/44}. \url{https://doi.org/10.4171/owr/2019/44}. This citation is supplied provenance; the report has not been independently consulted for this draft.
\item Supplied literature source: OWR 2019 44. \url{https://ems.press/content/serial-article-files/46824}.
\item Literature assessment checked at \texttt{2026-08-22}: supplied dated curation, not independently verified. It labels the problem open, describes the group-level realization as a natural question, mentions Bernstein projectives as a proposed route, and reports that no general construction or resolving theorem was found in that curation. This does not establish the absence of a solution.
\item The registry labels the statement \texttt{corrected\_verified} and says that notation, typography, and sentence structure were repaired while preserving mathematical content. These are supplied editorial claims, not mathematical verification of this draft.
\end{itemize}
\end{document}
